\begin{frame}{Pourquoi contraindre le XML?}
    \begin{itemize}
        \item Uniformiser l'encodage~;
        \item Collaborer à plusieurs sur un même projet~;
        \item Automatisation.
    \end{itemize}
\end{frame}

\begin{frame}{Spécifications du XML: \textit{Document Type Definition} (DTD)}
    La \textbf{DTD} est une \textbf{liste de déclarations} qui permet de définir:
    \begin{itemize}
        \item la structure d'un document~;
        \item les éléments et attributs autorisés.
    \end{itemize}
    \textcolor{red}{\textrightarrow{} Un document XML peut être bien formé mais invalide car ne respecte par les règles de la DTD !}
\end{frame}

\begin{frame}[fragile]{Spécifications du XML: \textit{Document Type Definition} (DTD)}
    \begin{block}{Exemple (source W3C):}
\begin{lstlisting}
<!DOCTYPE NEWSPAPER [

<!ELEMENT NEWSPAPER (ARTICLE+)>
<!ELEMENT ARTICLE (HEADLINE,BYLINE,LEAD,BODY,NOTES)>
<!ELEMENT HEADLINE (#PCDATA)>
<!ELEMENT BYLINE (#PCDATA)>
<!ELEMENT LEAD (#PCDATA)>
<!ELEMENT BODY (#PCDATA)>
<!ELEMENT NOTES (#PCDATA)>

<!ATTLIST ARTICLE AUTHOR CDATA #REQUIRED>
<!ATTLIST ARTICLE EDITOR CDATA #IMPLIED>
<!ATTLIST ARTICLE DATE CDATA #IMPLIED>
<!ATTLIST ARTICLE EDITION CDATA #IMPLIED>
\end{lstlisting}

    \end{block}
        
\end{frame}

\begin{frame}{Spécifications du XML: \textit{Document Type Definition} (DTD)}
    Limites du DTD:
    \begin{itemize}
        \item N'est pas rédigé en XML~;
        \item Assez limité en terme d'attributs~;
        \item Difficile à lire~;
        \item Limité en espace de noms.
    \end{itemize}
\end{frame}

\begin{frame}{Les schémas XML: définition}
\begin{block}{}
\textbf{Définition:}\\
    \small{
    \og Un schéma spécifie \textbf{un ensemble de noms d'éléments}, \textbf{les noms et le type de données de tous les attributs} qui leur sont associés, et \textbf{les règles relatives aux contextes} dans lesquels ils peuvent apparaître.\fg{}}\\
    Burnard, Lou. \og La TEI et le XML \fg. \textit{Qu'est-ce que la Text Encoding Initiative }?, OpenEdition Press, 2015, \url{https://doi.org/10.4000/books.oep.1298}.
\end{block}
    
\end{frame}

\begin{frame}{Les schémas XML: exemple du XML-EAD}
    \textbf{\textit{Encoded Archival Description} =} crée dans les années 1990, c'est une DTD utilisée pour la description de fonds d'archive. Depuis 2007, est un format XML.\\
     \textbf{Exemple:} \href{https://ccfr.bnf.fr/portailccfr/servlet/ViewManager?menu=public_menu_view&record=eadbam:EADC:d0e1129770_FRBNFEAD0000129462&setCache=all_simple.PUBLIC_CATALOGUE_SIMPLE_MULTI&fromList=true&source=all_simple&action=public_search_federated&query=PUBLIC_CATALOGUE_SIMPLE}{Dossiers individuels de prisonniers de la Bastille}
\end{frame}


